\documentclass[10pt]{article}
\usepackage[utf8]{inputenc}
\usepackage[vietnamese]{babel}
\usepackage{amsmath,amssymb,amsfonts}
\usepackage{graphicx}
\usepackage{xcolor}
\usepackage[most]{tcolorbox}
\usepackage{newtxtext,newtxmath}
\usepackage[paperwidth=216mm, paperheight=279mm, top=18mm, bottom=20mm, left=18mm, right=18mm, headheight=14pt, footskip=12mm]{geometry}
\usepackage{fancyhdr}

% Định nghĩa màu sắc nhận diện sách Stewart Calculus
\definecolor{stewartcyan}{RGB}{0, 136, 180}
\definecolor{stewartred}{RGB}{205, 35, 38}
\definecolor{stewartgray}{RGB}{100, 100, 100}

% Cấu hình Header & Footer
\pagestyle{fancy}
\fancyhf{}
\renewcommand{\headrulewidth}{0pt}
\fancyhead[L]{\fontsize{9}{11}\selectfont \textbf{130}\qquad \textbf{\textsf{CHƯƠNG 2}}\quad \textsf{Giới hạn và Đạo hàm}}
\fancyfoot[C]{%
    \begin{minipage}{\textwidth}
        \centering
        \tiny\color{stewartgray}
        Copyright 2021 Cengage Learning. All Rights Reserved. May not be copied, scanned, or duplicated, in whole or in part. WCN 02-200-203\\[1pt]
        Editorial review has deemed that any suppressed content does not materially affect the overall learning experience. Cengage Learning reserves the right to remove additional content at any time if subsequent rights restrictions require it.
    \end{minipage}
}

\setlength{\parindent}{0pt}
\setlength{\parskip}{5pt}

\begin{document}

\noindent
\begin{minipage}[t]{0.24\textwidth}
    % Cột lề trái: chứa Hình 7 dóng hàng chính xác với các bước tính của Ví dụ 3
    \vspace*{12.2cm}
    \centering
    \includegraphics[width=\linewidth]{images/p0165_fig1.png}
    
    \vspace{4pt}
    \begin{flushleft}
        \fontsize{9}{11}\selectfont
        \textbf{\textsf{HÌNH 7}}\\[2pt]
        $y = \dfrac{3x^2 - x - 2}{5x^2 + 4x + 1}$
    \end{flushleft}
\end{minipage}%
\hfill
\begin{minipage}[t]{0.72\textwidth}
    % Cột nội dung chính
    \vspace{0pt}
    {\color{stewartcyan}\rule{1.1ex}{1.1ex}}\;\; \textbf{\textsf{\large Tính giới hạn tại vô cực}}

    \vspace{2pt}
    Hầu hết các Quy tắc giới hạn đã được trình bày trong Mục 2.3 cũng nghiệm đúng đối với các giới hạn tại vô cực. Người ta có thể chứng minh rằng \textit{các Quy tắc giới hạn được liệt kê trong Mục 2.3 (ngoại trừ các Quy tắc 10 và 11) vẫn đúng nếu ``$x \to a$'' được thay thế bởi ``$x \to \infty$'' hoặc ``$x \to -\infty$''.} Đặc biệt, nếu ta kết hợp Quy tắc 6 và 7 với các kết quả của Ví dụ 2, ta thu được quy tắc quan trọng sau đây để tính giới hạn.

    \vspace{6pt}
    \begin{tcolorbox}[
        enhanced,
        colback=white,
        colframe=stewartred,
        boxrule=0.9pt,
        sharp corners,
        top=8pt, bottom=8pt, left=12pt, right=12pt
    ]
        \tcbox[
            on line,
            colback=white,
            colframe=stewartred,
            boxrule=0.8pt,
            sharp corners,
            boxsep=0pt,
            left=3pt, right=3pt, top=1.5pt, bottom=1.5pt
        ]{\textbf{\textsf{\color{stewartred}5}}}\;\;
        \textbf{\textsf{\color{stewartred}Định lý}}\quad
        Nếu $r > 0$ là một số hữu tỉ, thì
        \vspace{-2pt}
        \[
        \lim_{x \to \infty} \frac{1}{x^r} = 0
        \]
        \vspace{-6pt}
        Nếu $r > 0$ là một số hữu tỉ sao cho $x^r$ xác định với mọi $x$, thì
        \vspace{-2pt}
        \[
        \lim_{x \to -\infty} \frac{1}{x^r} = 0
        \]
        \vspace{-6pt}
    \end{tcolorbox}

    \vspace{6pt}
    \textbf{\textsf{\color{stewartcyan}VÍ DỤ 3}}\quad Tính giới hạn sau và chỉ rõ các tính chất của giới hạn được sử dụng ở mỗi bước.
    \[
    \lim_{x \to \infty} \frac{3x^2 - x - 2}{5x^2 + 4x + 1}
    \]
    \textbf{\textsf{\color{stewartcyan}LỜI GIẢI}}\quad Khi $x$ trở nên lớn, cả tử thức và mẫu thức đều trở nên lớn, do đó không thể thấy ngay tỉ số của chúng sẽ tiến tới đâu. Chúng ta cần thực hiện một số phép biến đổi đại số sơ bộ.

    Để tính giới hạn tại vô cực của một hàm phân thức hữu tỉ bất kỳ, trước hết ta chia cả tử thức và mẫu thức cho lũy thừa cao nhất của $x$ xuất hiện trong mẫu thức. (Ta có thể giả định rằng $x \neq 0$, vì ta chỉ quan tâm đến các giá trị lớn của $x$.) Trong trường hợp này, lũy thừa cao nhất của $x$ ở mẫu thức là $x^2$, vì vậy ta có
    \begin{align*}
    \lim_{x \to \infty} \frac{3x^2 - x - 2}{5x^2 + 4x + 1}
    &= \lim_{x \to \infty} \frac{\dfrac{3x^2 - x - 2}{x^2}}{\dfrac{5x^2 + 4x + 1}{x^2}} = \lim_{x \to \infty} \frac{3 - \dfrac{1}{x} - \dfrac{2}{x^2}}{5 + \dfrac{4}{x} + \dfrac{1}{x^2}} \\[8pt]
    &= \frac{\lim_{x \to \infty} \left( 3 - \dfrac{1}{x} - \dfrac{2}{x^2} \right)}{\lim_{x \to \infty} \left( 5 + \dfrac{4}{x} + \dfrac{1}{x^2} \right)} \tag*{\color{stewartcyan}\textsf{\footnotesize (theo Quy tắc giới hạn 5)}} \\[8pt]
    &= \frac{\lim_{x \to \infty} 3 - \lim_{x \to \infty} \dfrac{1}{x} - 2\lim_{x \to \infty} \dfrac{1}{x^2}}{\lim_{x \to \infty} 5 + 4\lim_{x \to \infty} \dfrac{1}{x} + \lim_{x \to \infty} \dfrac{1}{x^2}} \tag*{\color{stewartcyan}\textsf{\footnotesize (theo 1, 2, và 3)}} \\[8pt]
    &= \frac{3 - 0 - 0}{5 + 0 + 0} \tag*{\color{stewartcyan}\textsf{\footnotesize (theo 8 và Định lý 5)}} \\[8pt]
    &= \frac{3}{5}
    \end{align*}

    Một phép tính tương tự cho thấy giới hạn khi $x \to -\infty$ cũng bằng $\frac{3}{5}$. Hình 7 minh họa
\end{minipage}

\end{document}